% =====================================================================
%  ch08.tex —— 第 8 章 环面流形（原书 p.102–118 / PDF 111–127）
% =====================================================================
\chapter{环面流形}

% [原文 p.102]
$2n$ 维辛流形上的 Hamilton 系若具有 $n$ 个泊松可交换且函数独立的守恒量，便称为完全可积系统。由 Liouville--Arnold 定理，完全可积系统与 Lagrange 环面纤维化几乎是同一概念；但更强地，当 $2n$ 维紧辛流形具有 $n$ 维环面的有效 Hamilton 作用时，称之为\emph{辛环面流形}（symplectic toric manifold）。辛环面流形与满足适当条件的凸多面体一一对应，于是环面流形（toric manifold）的几何结构，可借助凸多面体的组合论结构以可见的形式加以理解。在这个意义上，可以说环面流形是凸体的组合论得以实体化的空间。环面流形在辛几何中发挥着超出单纯提供例子的重要作用。

% [Tutor Note: 「辛环面流形 $\leftrightarrow$ 凸多面体」的对应即下一节的 Delzant 定理；它把辛几何问题翻译成凸多面体的组合论问题，故常被称为辛几何的「离散化」。]

\section{辛环面流形}

对辛流形 $(X,\omega)$ 上的完全可积系统 $F\colon X\to\R^n$，称 $F$ 的正则值的原像为\emph{正则纤维}（regular fiber）。由 Liouville--Arnold 定理，若 $F$ 的正则纤维紧致且连通，则它成为 Lagrange 环面，并且在其邻域内可取到称为\emph{作用--角变量}的良好坐标。旋转角坐标使环面作用于 $F$ 的正则纤维，而作用坐标正是该作用的动量映射。作用--角坐标是这样一个 Darboux 坐标的类似物：沿纤维方向是大范围的，而沿底空间方向是局部的；它带有由 $GL_n(\Z)\ltimes\R^n$ 的仿射变换造成的不确定性。由于这种不确定性，作用--角坐标一般具有单值性（monodromy），不能大范围地定义。

% [Tutor Note: 单值性的来源是 $GL_n(\Z)\ltimes\R^n$ 中仿射部分与线性部分的非平凡粘合：绕底空间中的回路一周回到原纤维时，作用坐标被整体平移、角坐标被 $GL_n(\Z)$ 整变换。]

作用--角坐标没有单值性，且以作用坐标为 Hamilton 函数的 Hamilton 流给出 $n$ 维实环面 $T\cong\R^n/\Z^n$ 在 $X$ 上的作用，这个条件非常强，因此这种情形被赋予了下述名称：

% [原文 p.103]
\begin{definition}{8.1.1}
  紧致 $2n$ 维辛流形 $(X,\omega)$ 与 $n$ 维环面 $T\cong\R^n/\Z^n$ 在 $(X,\omega)$ 上的有效 Hamilton 作用这一组合 $(X,\omega,T,\mu)$，称为\emph{辛环面流形}（symplectic toric manifold）。
\end{definition}

这里，群 $G$ 在集合 $X$ 上的作用称为\emph{有效的}（effective），是指对任意 $g\in G$ 都存在某个 $x\in X$ 使得 $g\cdot x\neq x$。有效也称为\emph{忠实}（faithful）的。当 Lie 群 $G$ 作用于流形 $X$ 时，
\begin{equation}
  Z:=\{g\in G\mid \text{对任意 }x\in X,\ g\cdot x=x\} \tag{8.1.1}
\end{equation}
是 $G$ 的闭子群，因此若 $G$ 是可换群，则必要时把 $G$ 换为 $G/Z$，便可假定群作用总是有效的，须注意这一点。

紧致性的假设并非本质的，非紧的辛环面流形也可自然地定义。在复射影空间上赋予 Fubini--Study 形式所得者，是辛环面流形的典型例子；与此类似，在复仿射空间 $\C^n$ 上赋予由 Euclid 度量诱导的凯勒形式所得者，则是非紧辛环面流形的典型例子。此外，后述的三维普通二重点（ordinary double point）的 crepant 奇点消解，给出非紧辛环面流形的另一个典型例子。

% [Tutor Note: 「crepant 奇点消解」指消解映射 $f\colon\tilde X\to X$ 满足 $K_{\tilde X}=f^*K_X$（典范除子不改变）的奇点消解；三维普通二重点即 8.4.3 节 Atiyah flop 的出发点。]

具有半维数环面（half-dimensional torus）的有效 Hamilton 作用这一条件非常强，因此仅凭动量映射的像，就足以连同环面作用的信息一起决定其作为辛流形的结构。这称为 Delzant 定理，是辛环面流形的基本定理。为准确叙述它，先作些准备。

称有限生成自由阿贝尔群为\emph{格子}（lattice）。对 $n$ 维实环面 $T\cong\R^n/\Z^n$，用下式定义由单参数子群构成的格子 $N$ 与由指标构成的格子 $M$：
\begin{equation}
  N=\Hom(\R/\Z,T),\qquad M=\Hom(T,\R/\Z). \tag{8.1.2}
\end{equation}
其中 $N_\R:=N\otimes_\Z\R$ 与 $T$ 的李代数 $\mathfrak{t}$ 自然同构，且 $N$ 与 $M$ 互为对偶：
\begin{equation}
  M=\Hom(N,\Z),\qquad N=\Hom(M,\Z). \tag{8.1.3}
\end{equation}

称 $\R^n$ 中有限个点的凸包为\emph{多面体}（polytope）。

\begin{definition}{8.1.2}
  考虑秩 $n$ 的格子 $M$ 与 $M_\R$ 内的多面体 $\Delta$。
  \begin{enumerate}[label=(\arabic*),leftmargin=2.2em,itemsep=2pt]
    \item 对多面体 $\Delta$ 的顶点 $p$，若过 $p$ 的 $\Delta$ 的边恰有 $n$ 条，则称顶点 $p$ 是\emph{单纯的}（simple）。当所有顶点都单纯时，称多面体 $\Delta$ 是\emph{单纯的}（simple）。
% [原文 p.104]
    \item 当 $\Delta$ 的边 $e$ 的倾斜是\emph{有理的}时（即存在适当的 $p\in M\otimes\R$、$w\in M$ 及 $a,b\in\R$ 使
      \begin{equation}
        e=\{p+tw\mid t\in[a,b]\} \tag{8.1.4}
      \end{equation}
      时），称边 $e$ 是\emph{有理的}（rational）。当所有边都有理时，称多面体 $\Delta$ 是\emph{有理的}（rational）。
    \item
      单纯且有理的多面体 $\Delta$，在满足下述条件时称为\emph{光滑的}（smooth）：对任意顶点 $p$，把过 $p$ 的 $n$ 条边 $\{e_i\}_{i=1}^{n}$ 分别用适当的 $p_i\in M\otimes\R$、$w_i\in M$ 及 $a_i,b_i\in\R$ 表为
% [Tutor Note: 原书此条行文作「多面体の $\Delta$ の単純かつ有理的な多面体 $\Delta$ は」，文句重复，疑为「多面体 $\Delta$ 是单纯且有理的，且满足下述条件时称为光滑」之讹；此处照其义译出。]
      \begin{equation}
        e_i=\{p+tw_i\mid t\in[a_i,b_i]\} \tag{8.1.5}
      \end{equation}
      时，$\{w_i\}_{i=1}^{n}$ 构成 $M$ 的生成元。
  \end{enumerate}
\end{definition}

光滑多面体也称为\emph{Delzant 多面体}（Delzant polytope）。为固定平移的自由度，考虑 Delzant 多面体时总假定原点为重心。Delzant 多面体这一名称来自下述定理：

\begin{theorem}{8.1.3}（Delzant [33]）
  \begin{enumerate}[label=(\arabic*),leftmargin=2.2em,itemsep=2pt]
    \item 对任意辛环面流形，环面作用的动量映射之像都是 Delzant 多面体。
    \item 对任意 Delzant 多面体 $\Delta$，存在某个辛环面流形，其环面作用的动量映射之像与 $\Delta$ 一致。
    \item 两个辛环面流形作为辛环面流形同构的充要条件是：它们的动量映射之像作为多面体同构。
  \end{enumerate}
\end{theorem}

% [Tutor Note: 定理 8.1.3 是本章主干：它把「辛环面流形的同构分类」完全归结为「Delzant 多面体的同构分类」，即一个纯组合论问题；(2) 的构造在 p.105--106 给出。]

这里，对格子 $M$ 与 $M'$，称 $\Delta\subset M_\R$ 与 $\Delta'\subset M'_\R$ \emph{作为多面体同型}，是指存在格子的同构 $\phi\colon M\to M'$，使 $\phi$ 诱导的线性映射 $\phi_\R:=\phi\otimes\R\colon M_\R\to M'_\R$ 把 $\Delta$ 移到 $\Delta'$。又，两个辛环面流形 $(X,\omega,T,\mu)$ 与 $(X',\omega',T',\mu')$ 之间的\emph{同构映射}，是指 $n$ 维实环面 $T$ 与 $T'$ 之间作为 Lie 群的同构 $\psi\colon T\to T'$，以及 $(X,\omega)$ 与 $(X',\omega')$ 之间的辛同相 $\varphi\colon(X,\omega)\to(X',\omega')$，使得图式
\begin{equation}
  \begin{array}{ccc}
    X & \xrightarrow{\ \varphi\ } & X'\\[2pt]
    {\scriptstyle\mu}\big\downarrow & & \big\downarrow{\scriptstyle\mu'}\\[2pt]
    M_\R & \xleftarrow{\ \psi^{*}\ } & M'_\R
  \end{array}
  \tag{8.1.6}
\end{equation}
可交换。其中，把 $\psi$ 诱导的从 $M'_\R\cong(\mathfrak{t}')^*$ 到 $M_\R\cong\mathfrak{t}^*$ 的同构映射记作 $\psi^*$。

Hamilton 作用由动量映射唯一确定，注意到这一点，则图式 (8.1.6) 的可交换性等价于：在借助 $\psi$ 把两个环面 $T$ 与 $T'$ 等同起来时，辛同相 $\varphi$ 关于环面作用是等变的。以下沿 [30] 叙述 Delzant 定理的证明思路。

% [原文 p.105]
对环面流形 $(X,\omega,T,\mu)$ 的任意点 $x\in X$，存在 $0$ 与 $n$ 之间的整数 $r$，使 $T$ 的作用在 $x$ 处的稳定化子群成为 $r$ 维子环面 $T_x\subset T$。适当选取环面的坐标（即 Lie 群同构 $T\cong\R^n/\Z^n$），可使 $T_x$ 对应于 $\R^n/\Z^n$ 的子群
\begin{equation}
  \{(\alpha_1,\dots,\alpha_n)\in\R^n/\Z^n\mid \alpha_1=\cdots=\alpha_r=0\} \tag{8.1.7}
\end{equation}
此时，存在 $x$ 的开邻域 $U_x$，它与
\begin{equation}
  \left(\C^r\times(\C^\times)^{n-r},\ \frac{\sqrt{-1}}{2}\left(\sum_{i=1}^{r}\rd z_i\wedge\rd\bar z_i+\sum_{i=r+1}^{n}\frac{\rd z_i\wedge\rd\bar z_i}{|z_i|^2}\right)\right) \tag{8.1.8}
\end{equation}
的 $T$ 不变开集，作为带有 $T$ 的 Hamilton 作用的流形而同型。其中 $T$ 按
\begin{equation}
  (\alpha_1,\dots,\alpha_n)\colon (z_i)_{i=1}^{n}\mapsto\bigl(e^{2\pi\sqrt{-1}\alpha_i}z_i\bigr)_{i=1}^{n} \tag{8.1.9}
\end{equation}
作用于 $\C^r\times(\C^\times)^{n-r}$，且在此同型下 $x$ 移到 $(0,\dots,0,1,\dots,1)$。该作用的动量映射为
\begin{equation}
  \mu(z)=\Bigl(\tfrac12|z_1|^2,\dots,\tfrac12|z_r|^2,\log|z_{r+1}|,\dots,\log|z_n|\Bigr) \tag{8.1.10}
\end{equation}
它给出环面 Hamilton 作用的局部标准形。

% [Tutor Note: (8.1.8) 中第二项 $\rd z\wedge\rd\bar z/|z|^2$ 恰是 $\C^\times$ 上不变形式 $\rd\log|z|\wedge\rd\arg z$ 的两倍，故其动量映射取 $\log|z|$；第一项对应 $\C$ 上的标准形式，动量映射为 $|z|^2/2$。]

由上述局部标准形可知：对动量多面体 $\Delta$ 的任意点 $m$，存在 $m$ 的充分小邻域 $U$、格子 $M$ 与 $\Z^n$ 的适当等同，以及适当的 $\lambda_1,\dots,\lambda_r\in\R$，使得
\begin{equation}
  \Delta\cap U=\{m\in\R^n\mid \text{对 }i=1,\dots,r\text{ 均有 }x_i\ge\lambda_i\} \tag{8.1.11}
\end{equation}
这里 $\lambda_1,\dots,\lambda_r$ 的自由度来自动量映射定义中可以加常数的任意性。由这一表示立即可得：辛环面流形的动量多面体是 Delzant 多面体。

反之，给定 Delzant 多面体 $\Delta$ 时，以 $\Delta$ 为动量多面体的辛环面流形可按如下方式构造。首先，记 $\Delta$ 的余维数 $1$ 的面的集合为 $\Delta^{(1)}=\{F_1,\dots,F_d\}$，对每个 $F_i$，取其本原的内向法向量 $u_i\in N$，则 $\Delta$ 可由适当的 $\lambda_1,\dots,\lambda_d\in\R$ 表为
\begin{equation}
  \Delta=\{m\in M_\R\mid \text{对任意 }i=1,\dots,d,\ \langle u_i,m\rangle\ge\lambda_i\} \tag{8.1.12}
\end{equation}
% [原文 p.106]
用这些 $u_i$，按下式定义阿贝尔群的同态：
\begin{equation}
  \varphi\colon\Z^d\to N,\qquad \varphi(e_i)=u_i \tag{8.1.13}
\end{equation}
其中把 $\Z^d$ 的第 $i$ 个标准基记作 $e_i$。将 $\varphi$ 与 $\R/\Z$ 作张量积，得到环面的同态
\begin{equation}
  \varphi_{\R/\Z}\colon\R^d/\Z^d\to T=N_\R/N \tag{8.1.14}
\end{equation}
其核记为
\begin{equation}
  K:=\Ker\varphi_{\R/\Z}. \tag{8.1.15}
\end{equation}

$d$ 维实环面 $\R^d/\Z^d$ 具有在实 $2d$ 维辛向量空间 $(\C^d,\omega_{\C^d})$ 上的自然 Hamilton 作用。把该作用限制到 $\R^d/\Z^d$ 的子群 $K$，便得到 $K$ 在 $(\C^d,\omega_{\C^d})$ 上的 Hamilton 作用。该作用的动量映射 $\mu_K$，由 $\R^d/\Z^d$ 的动量映射 $\mu_{\R^d/\Z^d}$ 与嵌入映射 $\iota\colon K\hookrightarrow\R^d/\Z^d$ 的微分 $\iota_*\colon\Lie(\R^d/\Z^d)\to\mathfrak{k}$ 所诱导的拉回 $\iota^*\colon(\Lie(\R^d/\Z^d))^*\to\mathfrak{k}^*$ 合成给出；
\begin{equation}
  \mu_K=\iota^*\circ\mu_{\R^d/\Z^d}. \tag{8.1.16}
\end{equation}
此时，把 $\C^d$ 关于 $K$ 在 $\iota^*(-\lambda)$ 处的辛商记作
\begin{equation}
  X:=\C^d/\!/_{\iota^*(-\lambda)}K:=\mu_K^{-1}\bigl(\iota^*(-\lambda)\bigr)/K \tag{8.1.17}
\end{equation}

% [Tutor Note: 此处辛商是在 $\mu_K$ 的「平移」水平集上取的；$\iota^*(-\lambda)$ 是 $\mathfrak{k}^*$ 中的正则值，Delzant 多面体的法向量数据保证 $K$ 作用自由，故商空间光滑。]

Lie 群的短正合列
\begin{equation}
  1\to K\xrightarrow{\ \iota\ }\R^d/\Z^d\xrightarrow{\ \varphi\otimes(\R/\Z)\ }N_\R/N\to 1 \tag{8.1.18}
\end{equation}
所诱导的李代数的短正合列
\begin{equation}
  0\to\mathfrak{k}\xrightarrow{\ \iota_*\ }\R^d\xrightarrow{\ \varphi_*\ }N_\R\to 0 \tag{8.1.19}
\end{equation}
称为\emph{扇列}（fan sequence），而取其对偶所得的实向量空间的短正合列
\begin{equation}
  0\to M_\R\xrightarrow{\ \varphi^*\ }\R^d\xrightarrow{\ \iota^*\ }\mathfrak{k}^*\to 0 \tag{8.1.20}
\end{equation}
称为\emph{除子列}（divisor sequence）。$X$ 具有 $N_\R/N$ 的有效 Hamilton 作用，其动量多面体由 $\R^d$ 的第一象限
\begin{equation}
  \{(x_1,\dots,x_d)\in\R^d\mid \text{对任意 }i=1,\dots,d,\ x_i\ge 0\} \tag{8.1.21}
\end{equation}
与 $(\iota^*)^{-1}\bigl(\iota^*(-\lambda)\bigr)$ 的公共部分给出。以点 $-\lambda=-(\lambda_1,\dots,\lambda_d)\in\R^d$ 为原点，把仿射空间 $(\iota^*)^{-1}\bigl(\iota^*(-\lambda)\bigr)$ 与向量空间 $M_\R$ 等同起来，则该公共部分与 (8.1.12) 所定义的 $\Delta$ 一致。

% [Tutor Note: 扇列与除子列的对偶关系是 Delzant 构造的枢纽：组合数据（多面体的法向量 $u_i$）经 $\varphi$ 变成环面同态，其核 $K$ 再经辛商还原出 $X$；除子列则说明 $M_\R$ 正是 $\R^d$ 中被「除掉」的 $K$ 方向的正交补。]

\takeaway{半维数环面的有效作用使辛流形由 Delzant 多面体完全决定，二者一一对应。}

% [原文 p.107]
要证明辛环面流形由动量多面体唯一决定，只需对共享动量多面体 $\Delta$ 的任意两个辛环面流形 $X$ 与 $X'$，构造 $X$ 与 $X'$ 之间的同型映射即可。为此，首先把前述的局部标准形拼接起来，构造从 $X$ 到 $X'$ 的微分同相，它关于群作用是等变的，并保持动量映射。借助该微分同相把 $X$ 与 $X'$ 等同起来，便得到 $X$ 上两个环面不变的辛形式 $\omega$ 与 $\omega'$，它们对应的动量映射作为映射相等。此时，对任意 $t\in[0,1]$，
\begin{equation}
  \omega_t=(1-t)\omega+t\omega' \tag{8.1.22}
\end{equation}
仍是环面不变的辛形式。容易看出，环面不变辛形式所属的上同调类由动量多面体决定，因此 $\omega_t$ 的上同调类不依赖于 $t$。这里使用 Moser 定理，最终可得 $X$ 与 $X'$ 作为辛环面流形同型。

% [Tutor Note: 用 Moser 定理时，关键在于 $\omega_t$ 的上同调类恒定，故 $\dot\omega_t=\rd\sigma_t$ 有原像，解 Moser 方程 $\iota_{v_t}\omega_t=-\sigma_t$ 即得一族保持 $[\omega_t]$ 的微分同相。]

\section{作为代数簇的环面簇}

所有辛环面流形都自然具有复射影簇的结构，因此也可以从复代数几何的观点研究这些流形。在复代数几何中，环面簇按下述方式定义：

\begin{definition}{8.2.1}
  若代数的环面 $T_\C=(\C^\times)^n$ 所作用的正规代数簇 $X$ 具有稠密的 $T_\C$ 轨道，则称 $X$ 连同其上的 $T_\C$ 作用这一组合为\emph{环面簇}（toric variety）。
\end{definition}

注意这里用的是 variety（簇）而非 manifold（流形）。环面簇未必光滑，一般具有奇点。辛环面流形局部上与 (8.1.8) 同型，但环面簇的局部理论内容要丰富得多。

环面簇由与动量多面体相对偶的、称为\emph{扇形}的对象来描述。

\takeaway{环面簇由扇形把仿射概形 $\operatorname{Spec}\C[\sigma^\vee\cap M]$ 粘贴而成，全部组合数据集中在扇形。}

\begin{definition}{8.2.2}
  在 $\R^{\ge 0}$ 的作用下不变的 $N_\R$ 的子集称为\emph{锥}（cone）。对锥 $\sigma$，置
  \begin{equation}
    -\sigma:=\{-n\in N_\R\mid n\in\sigma\}. \tag{8.2.1}
  \end{equation}
  当锥 $\sigma$ 是凸集，即满足
  \begin{equation}
    \text{对任意 } n,n'\in\sigma \text{ 与 } t\in[0,1],\quad (1-t)n+tn'\in\sigma \tag{8.2.2}
  \end{equation}
  时，称之为\emph{凸的}（convex）。

  % [原文 p.108]
  当锥 $\sigma$ 满足
  \begin{equation}
    \sigma\cap(-\sigma)=\{0\} \tag{8.2.3}
  \end{equation}
  时，称之为\emph{强凸的}（strongly convex）。由有限个线性不等式定义的锥称为\emph{多面锥}（polyhedral cone）。当锥 $\sigma$ 与 $\sigma\cap N$ 的凸包一致，即满足
  \begin{equation}
    \sigma=\operatorname{Conv}(\sigma\cap N) \tag{8.2.4}
  \end{equation}
  时，称之为\emph{有理的}（rational）。这与 $\sigma$ 可用有限个 $m_1,\dots,m_d\in M$ 表成
  \begin{equation}
    \sigma=\{n\in N_\R\mid \langle m_i,n\rangle\ge 0,\ i=1,\dots,d\} \tag{8.2.5}
  \end{equation}
  是等价的。在 (8.2.5) 中把若干个不等号换成等号所得的集合，称为 $\sigma$ 的\emph{面}（face）。当 $\tau$ 是 $\sigma$ 的面时，写作 $\tau\prec\sigma$。有理多面锥 $\sigma$ 的面是 $\sigma$ 的边界的子集。零维的面称为\emph{顶点}（vertex），一维的面称为\emph{棱}（edge），余维数为 $1$ 的面称为\emph{切子面}（facet）。当有理强凸多面锥的集合 $\Sigma$ 满足以下两个条件时，称为\emph{扇形}（fan）：
  \begin{itemize}
    \item 对任意 $\sigma\in\Sigma$，$\sigma$ 的任意面仍是 $\Sigma$ 的元素；
    \item 对任意 $\sigma,\tau\in\Sigma$，它们的公共部分 $\sigma\cap\tau$ 仍是 $\Sigma$ 的元素。
  \end{itemize}
\end{definition}

对锥 $\sigma$，按下式定义其\emph{对偶锥}（dual cone）$\sigma^\vee\subset M_\R$：
\begin{equation}
  \sigma^\vee:=\{m\in M_\R\mid \text{对任意 } n\in\sigma,\ \langle n,m\rangle\ge 0\}. \tag{8.2.6}
\end{equation}
若锥 $\sigma$ 是有理多面锥，则 $\sigma^\vee\cap M$ 是有限生成半群。把由对应的半群环 $\C[\sigma^\vee\cap M]$ 所得的仿射概形写作
\begin{equation}
  U_\sigma:=\operatorname{Spec}\C[\sigma^\vee\cap M]. \tag{8.2.7}
\end{equation}

由于格子 $M$ 与代数环面 $T_\C$ 的指标群 $\Hom(T_\C,\C^\times)$ 自然等同，故与半群环 $\C[\sigma^\vee\cap M]$ 对应的仿射概形 $U_\sigma$ 上自然有 $T_\C$ 作用。

若有理多面锥 $\sigma$ 强凸，则 $\sigma^\vee\cap M$ 作为群生成 $M$，因此可知 $U_\sigma$ 以开子概形的形式包含 $\operatorname{Spec}\C[M]\cong T_\C$。

由对偶锥的定义，对任意 $\tau\prec\sigma$ 有 $\tau^\vee\supset\sigma^\vee$，这引起半群环的包含关系
\begin{equation}
  \C[\tau^\vee\cap M]\supset\C[\sigma^\vee\cap M]. \tag{8.2.8}
\end{equation}
它所引起的仿射概形的开嵌入

% [原文 p.109]
\begin{equation}
  U_\tau:=\operatorname{Spec}\C[\tau^\vee\cap M]\hookrightarrow U_\sigma:=\operatorname{Spec}\C[\sigma^\vee\cap M] \tag{8.2.9}
\end{equation}
把 $(U_\sigma)_{\sigma\in\Sigma}$ 粘贴起来所得的代数簇，称为扇形 $\Sigma$ 所伴随的\emph{环面簇}，记作 $X_\Sigma$：
\begin{equation}
  X_\Sigma=\bigcup_{\sigma\in\Sigma}U_\sigma. \tag{8.2.10}
\end{equation}

由于仿射概形的嵌入 (8.2.9) 是由半群的嵌入 $\sigma^\vee\subset\tau^\vee$ 引起的，故嵌入 (8.2.9) 关于 $T_\C$ 作用是等变的；于是各 $U_\sigma$ 上的 $T_\C$ 作用粘贴起来，定义出 $X_\Sigma$ 上的 $T_\C$ 作用。对任意 $\sigma\in\Sigma$，因 $\sigma^\vee\subset M$，故
\begin{equation}
  U_\sigma=\operatorname{Spec}\C[\sigma^\vee\cap M]\supset T_\C=\operatorname{Spec}\C[M], \tag{8.2.11}
\end{equation}
并且把嵌入 (8.2.9) 限制到 $\operatorname{Spec}\C[M]$ 上便成为恒等映射。由于 $T_\C$ 是 $U_\sigma$ 的稠密开集，$X_\Sigma$ 把 $T_\C$ 作为稠密开集包含在内。

当格子 $M$ 的子幺半群 $S$ 满足
\begin{itemize}
  \item 对任意 $x\in M$ 与 $k\in\Z^{>0}$，若 $kx\in S$ 则 $x\in S$，
\end{itemize}
这一条件时，称 $S$ 是\emph{饱和的}（saturated）。对一般的可换幺半群，只需把 $M$ 换成 $S$ 的 Grothendieck 群
\begin{equation}
  S^{\mathrm{group}}=S\times S/\!\sim \tag{8.2.12}
\end{equation}
即可。其中等价关系定义为：当存在某个 $z\in S$ 使 $x+y'+z=x'+y+z$ 时，$(x,y)\sim(x',y')$。当到 Grothendieck 群的标准映射是单射时，称该幺半群是\emph{整的}（integral）。整幺半群 $S$ 的半群环为正规（normal）的充要条件是 $S$ 饱和。幺半群 $\sigma^\vee\cap M$ 显然饱和，故对任意扇形 $\Sigma$，$X_\Sigma$ 是正规簇。

对动量多面体的点 $x\in\Delta$，按下式定义 $\Delta$ 在 $x$ 处的\emph{切锥}（tangent cone）：
\begin{equation}
  T_x\Delta=\{m\in M_\R\mid \text{存在正数 }t\text{ 使 }x+tm\in\Delta\}. \tag{8.2.13}
\end{equation}
切锥 $T_x\Delta\subset M_\R$ 的对偶锥称为\emph{法锥}（normal cone），记作 $N_x\Delta\subset N_\R$。当 $x$ 与 $x'$ 是 $\Delta$ 的同一个面的相对内点时，$\Delta$ 在 $x$ 处的法锥与在 $x'$ 处的法锥一致。称 $\Delta$ 的法锥的集合 $\bigcup_{x\in\Delta}\{T_x\Delta\}$ 为 $\Delta$ 的\emph{法扇形}（normal fan）。

% [Tutor Note: 此处集合应为 $\bigcup_{x\in\Delta}\{N_x\Delta\}$；原书印作 $T_x\Delta$，与「法锥的集合」不合，属笔误。]

对辛环面流形，记其动量多面体为 $\Delta$，则由 $\Delta$ 的法扇形所定的环面簇与原来的辛环面流形微分同相。在不变动动量多面体的组合论结构的前提下平行移动切子面，法扇形并不改变，因此动量多面体比扇形含有更多的信息。这些多余的信息描述着辛形式所属的上同调类。另一方面，非射影的环面簇不具有辛环面流形的结构，因此要描述所有的环面簇必须使用扇形。此外，可以在未必是复数域的域上定义，也是环面簇相对于辛环面流形的优点。

\section{反射多面体与 Fano 多面体}

% [原文 p.110]
称扇形的一维锥的本原生成元的凸包为\emph{扇多面体}（fan polytope）。当凸多面体 $\Delta\subset M_\R$ 的所有顶点都含于 $M$ 时，称之为\emph{凸格多面体}（convex lattice polytope）。当凸格多面体 $\Delta\subset N_\R$ 满足
\begin{itemize}
  \item 原点含于内部，即 $0\in\operatorname{Int}\Delta$；
  \item 对任意余维数为 $1$ 的面 $F\subset\Delta$，$F$ 的顶点集合构成 $N$ 的 $\Z$ 基，
\end{itemize}
这两个条件时，称之为\emph{光滑 Fano 多面体}（smooth Fano polytope）。反典范除子（anti-canonical divisor）为丰富（ample）的射影簇（projective manifold）称为 \emph{Fano 流形}（Fano manifold）。环面簇为 Fano 流形的充要条件是其扇多面体为光滑 Fano 多面体。

\takeaway{光滑 Fano 多面体与反射多面体以极对偶相连，给出 Fano 环面簇的组合分类数据。}

多面体 $\Delta\subset M$ 的\emph{极对偶}（polar dual）定义为
\begin{equation}
  \Delta^\circ=\{n\in N\mid \text{对任意 } m\in\Delta,\ \langle n,m\rangle\ge -1\}. \tag{8.3.1}
\end{equation}
若 $\Delta$ 把原点含于内部，则容易看出
\begin{equation}
  (\Delta^\circ)^\circ=\Delta. \tag{8.3.2}
\end{equation}
对格多面体 $\Delta$，当极对偶 $\Delta^\circ$ 再次成为格多面体时，称 $\Delta$ 是\emph{反射的}（reflexive）。显然反射多面体的极对偶仍是反射的。此外容易看出，反射多面体以原点为唯一的内点。

容易看出，光滑 Fano 多面体在一维只有 $1$ 个，二维只有 $5$ 个。三维光滑 Fano 多面体由 Batyrev [16] 与渡辺–渡邊 [136] 独立分类，共有 $18$ 个。Batyrev [17] 得到了由四维光滑 Fano 多面体构成的 $123$ 个的列表，但实际上遗漏了一个；佐藤 [120] 指出四维光滑 Fano 多面体共有 $124$ 个。五维光滑 Fano 多面体由 Kreuzer–Nill [81] 分类，六维与七维光滑 Fano 多面体由 Øbro [112] 分类，已知分别有 $866$ 个、$7622$ 个与 $72256$ 个。

Hensley [61] 证明了：当 $n$ 维格多面体 $\Delta$ 内部含有 $n\ge 1$ 个格点时，$\Delta$ 的体积可由仅依赖于 $n$ 与 $k$ 的常数从上方控制。此外，Lagarias 与 Ziegler [84] 证明了，任意 $n$ 维格多面体 $\Delta$ 都经适当的整仿射变换

% [Tutor Note: 此处「$n\ge 1$ 个格点」应为「$k\ge 1$ 个格点」；原书后文随即出现「$n$ 与 $k$ 的常数」，可证此为笔误。]

% [原文 p.111]
变换后含于一个 $n$ 维格立方体中，该立方体的边长为 $n\cdot n!\operatorname{vol}(\Delta)$。由这些事实立即得出：以原点为唯一内点的格多面体的同型类在每个维数下都是有限的。容易看出，一维与二维的反射多面体分别有 $1$ 个与 $16$ 个。三维与四维的反射多面体分别有 $4319$ 个与 $473{,}800{,}776$ 个，这由 Kreuzer 与 Skarke 给出。关于这一点，参见 [82] 及其参考文献。此外，Haase 与 Melnikov [58] 指出，任意格多面体都可成为某个足够高维数的反射多面体的面。

反射多面体与其极对偶在镜像对称的研究中起着重要作用。关于这一话题，例如参见 [31] 及其引用文献。

\section{具体例}

\subsection{射影直线}

当 $\Delta$ 为一维时，令 $M=\Z$，把 $\Delta\subset M_\R$ 的定义不等式取为
\begin{equation}
  \Delta=\{m\in M_\R\mid m\ge\lambda_1\ \text{且}\ -m\ge\lambda_2\}, \tag{8.4.1}
\end{equation}
则阿贝尔群的同态 (8.1.13) 由下式给出：
\begin{equation}
  \varphi\colon\Z^2\to N\cong\Z,\qquad \varphi(1,0)=1,\quad \varphi(0,1)=-1. \tag{8.4.2}
\end{equation}
并且
\begin{align}
  K&:=\Ker\varphi_{\R/\Z} \tag{8.4.3}\\
   &=\{(\alpha_1,\alpha_2)\in\R^2/\Z^2\mid \alpha_1=\alpha_2\} \tag{8.4.4}\\
   &\cong S^1=\{\alpha\in\C^\times\mid |\alpha|=1\} \tag{8.4.5}
\end{align}
作用于 $\C^2$ 上：
\begin{equation}
  \alpha\colon(z_1,z_2)\mapsto(\alpha z_1,\alpha z_2). \tag{8.4.6}
\end{equation}
该作用是具有动量映射
\begin{equation}
  \mu_K\colon(z_1,z_2)\mapsto \frac{1}{2}\bigl(|z_1|^2+|z_2|^2\bigr) \tag{8.4.7}
\end{equation}
的 Hamilton 作用，对应的辛商为
\begin{equation}
  \C^2/\!/_\lambda K:=\mu_K^{-1}(\lambda)/K \tag{8.4.8}
\end{equation}
\begin{equation}
  \cong
  \begin{cases}
    \PP^1 & \lambda>0,\\
    1\ \text{点} & \lambda=0,\\
    \varnothing & \lambda<0.
  \end{cases} \tag{8.4.9}
\end{equation}

% [原文 p.112]
因此，扇列 (8.1.19) 成为
\begin{equation}
  0\to\R\xrightarrow{\ \iota_*=\left(\begin{smallmatrix}1\\)\ }\R^2\xrightarrow{\ \varphi_*=(1\ \ -1)\ }\R\to 0, \tag{8.4.10}
\end{equation}
因子列 (8.1.20) 成为
\begin{equation}
  0\to\R\xrightarrow{\ \varphi^*=\left(\begin{smallmatrix}1\\)\ }\R^2\xrightarrow{\ \iota^*=(1\ \ 1)\ }\R\to 0. \tag{8.4.11}
\end{equation}
对 $(-\lambda_1,-\lambda_2)\in\R^2$，有
\begin{equation}
  \iota^*(-\lambda_1,-\lambda_2)=-\lambda_1-\lambda_2, \tag{8.4.12}
\end{equation}
且
\[
  (\iota^*)^{-1}\bigl(\iota^*(-\lambda_1,-\lambda_2)\bigr)=\{(x_1,x_2)\in\R^2\mid x_1+x_2=-\lambda_1-\lambda_2\}.
\]
通过映射
\begin{equation}
  M_\R\to(\iota^*)^{-1}\bigl(\iota^*(-\lambda_1,-\lambda_2)\bigr),\qquad m\mapsto(m-\lambda_1,-m-\lambda_2) \tag{8.4.13}
\end{equation}
把这个仿射空间与 $M_\R\cong\R$ 等同起来，则 $(\iota^*)^{-1}\bigl(\iota^*(-\lambda_1,-\lambda_2)\bigr)$ 与第一象限 $\{(x_1,x_2)\in\R^2\mid x_1\ge 0,\ x_2\ge 0\}$ 的公共部分对应于 $\Delta\subset M_\R$。

$\Delta$ 的法扇形由 $2$ 个一维锥与 $1$ 个零维锥构成：
\begin{align}
  \Sigma&=\{\sigma_1,\sigma_2,\sigma_0\}, \tag{8.4.14}\\
  \sigma_0&=\{0\}, \tag{8.4.15}\\
  \sigma_1&=\{n\in N_\R\mid n\ge 0\}, \tag{8.4.16}\\
  \sigma_2&=\{n\in N_\R\mid n\le 0\}. \tag{8.4.17}
\end{align}
对应的 $M$ 的部分幺半群为
\begin{align}
  \sigma_0^\vee\cap M&=M, \tag{8.4.18}\\
  \sigma_1^\vee\cap M&=\{m\in M\mid m\ge 0\}, \tag{8.4.19}\\
  \sigma_2^\vee\cap M&=\{m\in M\mid m\le 0\}, \tag{8.4.20}
\end{align}
而半群环为
\begin{align}
  \C[\sigma_0^\vee\cap M]&\cong\C[t,t^{-1}], \tag{8.4.21}\\
  \C[\sigma_1^\vee\cap M]&\cong\C[t], \tag{8.4.22}\\
  \C[\sigma_2^\vee\cap M]&\cong\C[t^{-1}]. \tag{8.4.23}
\end{align}

% [原文 p.113]
由上式给出。与 $\sigma_1$ 及 $\sigma_2$ 对应的仿射概形 $U_{\sigma_1}$ 与 $U_{\sigma_2}$ 都与复平面 $\C$ 同型，把它们的自然坐标取为 $z=t$、$w=t^{-1}$，则 $X_\Sigma$ 是把 $U_{\sigma_1}$ 与 $U_{\sigma_2}$ 在
\begin{equation}
  U_{\sigma_0}=U_{\sigma_1}\cap U_{\sigma_2}\cong\C^\times \tag{8.4.24}
\end{equation}
处经变量变换
\begin{equation}
  w=z^{-1} \tag{8.4.25}
\end{equation}
粘贴起来所得的。

\takeaway{一个区间对应两个仿射平面经 $w=z^{-1}$ 的粘合，Delzant 构造由此还原射影直线。}

\subsection{射影平面}

其次考虑 $M=\Z^2$ 为二维、$\Delta$ 为由 $3$ 个不等式
\begin{align}
  m_1&\ge\lambda_1, \tag{8.4.26}\\
  m_2&\ge\lambda_2, \tag{8.4.27}\\
  -m_1-m_2&\ge\lambda_3 \tag{8.4.28}
\end{align}
定义的三角形的情形。此时，(8.1.13) 的阿贝尔群同态 $\varphi\colon\Z^3\to N\cong\Z^2$ 由
\[
  \varphi(1,0,0)=(1,0),\qquad \varphi(0,1,0)=(0,1),\qquad \varphi(0,0,1)=(-1,-1)
\]
给出，并且
\begin{equation}
  K=\{(\alpha_1,\alpha_2,\alpha_3)\in\R^3/\Z^3\mid \alpha_1-\alpha_3=\alpha_2-\alpha_3=0\}\cong S^1 \tag{8.4.29}
\end{equation}
作用于 $\C^3$ 上：
\begin{equation}
  \alpha\colon(z_1,z_2,z_3)\mapsto(\alpha z_1,\alpha z_2,\alpha z_3). \tag{8.4.30}
\end{equation}
该作用是具有动量映射
\begin{equation}
  \mu_K\colon(z_1,z_2,z_3)\mapsto\frac{1}{2}\bigl(|z_1|^2+|z_2|^2+|z_3|^2\bigr) \tag{8.4.31}
\end{equation}
的 Hamilton 作用，对应的辛商为
\begin{equation}
  \C^2/\!/_\lambda K:=\mu_K^{-1}(\lambda)/K \tag{8.4.32}
\end{equation}
\begin{equation}
  \cong
  \begin{cases}
    \PP^2 & \lambda>0,\\
    1\ \text{点} & \lambda=0,\\
    \varnothing & \lambda<0
  \end{cases} \tag{8.4.33}
\end{equation}

% [Tutor Note: (8.4.32) 左端原作 $\C^2$，但此处动量映射定义在 $\C^3$ 上、商为 $\PP^2$，故应为 $\C^3$；原书笔误。]

% [原文 p.114]
扇列 (8.1.19) 与因子列 (8.1.20) 分别由
\begin{equation}
  0\to\R\xrightarrow{\ \iota_*=\left(\begin{smallmatrix}1\\)\ }\R^3
  \xrightarrow{\ \varphi_*=\left(\begin{smallmatrix}1&0&-1\\0&1&-1\end{smallmatrix}\right)\ }\R^2\to 0 \tag{8.4.34}
\end{equation}
以及
\begin{equation}
  0\to\R^2\xrightarrow{\ \varphi^*=\left(\begin{smallmatrix}1&0\\0&1\\-1&-1\end{smallmatrix}\right)\ }\R^3
  \xrightarrow{\ \iota^*=(1\ \ 1\ \ 1)\ }\R\to 0 \tag{8.4.35}
\end{equation}
给出。并且
\[
  (\iota^*)^{-1}\bigl(\iota^*(-\lambda_1,-\lambda_2,-\lambda_3)\bigr)
  =\{(x_1,x_2,x_3)\in\R^3\mid x_1+x_2+x_3=-\lambda_1-\lambda_2-\lambda_3\},
\]
通过映射
\begin{align}
  M_\R&\to(\iota^*)^{-1}\bigl(\iota^*(-\lambda_1,-\lambda_2,-\lambda_3)\bigr), \tag{8.4.36}\\
  (m_1,m_2)&\mapsto(m_1-\lambda_1,\ m_2-\lambda_2,\ -m_1-m_2-\lambda_3) \tag{8.4.37}
\end{align}
把这个仿射空间与 $M_\R\cong\R^2$ 等同起来，则 $(\iota^*)^{-1}\bigl(\iota^*(-\lambda_1,-\lambda_2,-\lambda_3)\bigr)$ 与第一象限 $\{(x_1,x_2,x_3)\in\R^3\mid x_1\ge 0,\ x_2\ge 0,\ x_3\ge 0\}$ 的公共部分对应于 $\Delta\subset M_\R$。对应的扇形由 $3$ 个二维锥
\begin{align}
  \sigma_1&=\{(n_1,n_2)\in N_\R\mid n_1\ge 0,\ n_2\ge 0\}, \tag{8.4.38}\\
  \sigma_2&=\{(n_1,n_2)\in N_\R\mid -n_1\ge 0,\ -n_1+n_2\ge 0\}, \tag{8.4.39}\\
  \sigma_3&=\{(n_1,n_2)\in N_\R\mid n_1-n_2\ge 0,\ -n_2\ge 0\}, \tag{8.4.40}
\end{align}
$3$ 个一维锥
\begin{align}
  \sigma_{12}&=\sigma_1\cap\sigma_2=\R^{\ge 0}\cdot(0,1), \tag{8.4.41}\\
  \sigma_{23}&=\sigma_2\cap\sigma_3=\R^{\ge 0}\cdot(-1,-1), \tag{8.4.42}\\
  \sigma_{31}&=\sigma_3\cap\sigma_1=\R^{\ge 0}\cdot(1,0), \tag{8.4.43}
\end{align}
以及 $1$ 个零维锥
\begin{equation}
  \sigma_{123}=\sigma_1\cap\sigma_2\cap\sigma_3=\{0\} \tag{8.4.44}
\end{equation}
构成。对应的 $M$ 的部分幺半群为

% [原文 p.115]
\begin{align}
  \sigma_1^\vee\cap M&=\{(m_1,m_2)\in M\mid m_1\ge 0,\ m_2\ge 0\}, \tag{8.4.45}\\
  \sigma_2^\vee\cap M&=\{(m_1,m_2)\in M\mid m_2\ge 0,\ -m_1-m_2\ge 0\}, \tag{8.4.46}\\
  \sigma_3^\vee\cap M&=\{(m_1,m_2)\in M\mid m_1\ge 0,\ -m_1-m_2\ge 0\}, \tag{8.4.47}\\
  \sigma_{12}^\vee\cap M&=\{(m_1,m_2)\in M\mid m_2\ge 0\}, \tag{8.4.48}\\
  \sigma_{23}^\vee\cap M&=\{(m_1,m_2)\in M\mid -m_1-m_2\ge 0\}, \tag{8.4.49}\\
  \sigma_{31}^\vee\cap M&=\{(m_1,m_2)\in M\mid m_1\ge 0\}, \tag{8.4.50}\\
  \sigma_{123}^\vee\cap M&=M, \tag{8.4.51}
\end{align}
而半群环为
\begin{align}
  \C[\sigma_1^\vee\cap M]&\cong\C[t_1,t_2], \tag{8.4.52}\\
  \C[\sigma_2^\vee\cap M]&\cong\C[t_1^{-1},t_1^{-1}t_2], \tag{8.4.53}\\
  \C[\sigma_3^\vee\cap M]&\cong\C[t_1t_2^{-1},t_2^{-1}], \tag{8.4.54}\\
  \C[\sigma_{12}^\vee\cap M]&\cong\C[t_1,t_1^{-1},t_2], \tag{8.4.55}\\
  \C[\sigma_{23}^\vee\cap M]&\cong\C[t_1^{-1}t_2,t_1t_2^{-1},t_1^{-1}], \tag{8.4.56}\\
  \C[\sigma_{31}^\vee\cap M]&\cong\C[t_1,t_2,t_2^{-1}], \tag{8.4.57}\\
  \C[\sigma_{123}^\vee\cap M]&\cong\C[t_1,t_1^{-1},t_2,t_2^{-1}]. \tag{8.4.58}
\end{align}
对 $i=1,2,3$，取 $U_{\sigma_i}$ 的坐标为
\begin{align}
  U_{\sigma_1}:\quad t_1&=\frac{x}{z}, & t_2&=\frac{y}{z}, \tag{8.4.59}\\
  U_{\sigma_2}:\quad t_1^{-1}&=\frac{z}{x}, & t_1^{-1}t_2&=\frac{y}{x}, \tag{8.4.60}\\
  U_{\sigma_3}:\quad t_1t_2^{-1}&=\frac{x}{y}, & t_2^{-1}&=\frac{z}{y}, \tag{8.4.61}
\end{align}
则例如在 $U_{\sigma_{12}}=U_{\sigma_1}\cap U_{\sigma_2}$ 上
\begin{align}
  \frac{z}{x}&=t_1^{-1}=\Bigl(\frac{x}{z}\Bigr)^{-1}, \tag{8.4.62}\\
  \frac{y}{x}&=t_1^{-1}t_2=\Bigl(\frac{x}{z}\Bigr)^{-1}\frac{y}{z} \tag{8.4.63}
\end{align}
等等成立，在等同
\begin{align}
  U_{\sigma_1}&=\{[x:y:z]\in\PP^2\mid z\ne 0\}, \tag{8.4.64}\\
  U_{\sigma_2}&=\{[x:y:z]\in\PP^2\mid x\ne 0\}, \tag{8.4.65}\\
  U_{\sigma_3}&=\{[x:y:z]\in\PP^2\mid y\ne 0\} \tag{8.4.66}
\end{align}
之下，可知
\begin{equation}
  X_\Sigma=U_{\sigma_1}\cup U_{\sigma_2}\cup U_{\sigma_3}=\PP^2 \tag{8.4.67}
\end{equation}
成立。

\subsection{Atiyah flop}

% [原文 p.116]
一维锥由
\[
  \sigma_1=\R^{\ge 0}\cdot(0,0,1),\quad \sigma_2=\R^{\ge 0}\cdot(1,1,1),\quad
  \sigma_3=\R^{\ge 0}\cdot(1,0,1),\quad \sigma_4=\R^{\ge 0}\cdot(0,1,1)
\]
给出的非紧辛环面流形有两个，它们通过称为 \emph{Atiyah flop}（Atiyah flop）的操作相联系。此时的扇列为
\begin{equation}
  0\to\R\xrightarrow{\ (1\ \ 1\ \ -1\ \ -1)\ }\R^4
  \xrightarrow{\ \left(\begin{smallmatrix}0&1&1&0\\0&1&0&1\\1&1&1&1\end{smallmatrix}\right)\ }\R^3\to 0 \tag{8.4.68}
\end{equation}
$K\cong S^1$ 在 $\C^4$ 上的 Hamilton 作用为
\begin{equation}
  S^1\ni\alpha\colon(z_1,z_2,z_3,z_4)\mapsto(\alpha z_1,\alpha z_2,\alpha^{-1}z_3,\alpha^{-1}z_4), \tag{8.4.69}
\end{equation}
其动量映射为
\begin{equation}
  \mu_K(z_1,z_2,z_3,z_4)=\frac{1}{2}\bigl(|z_1|^2+|z_2|^2-|z_3|^2-|z_4|^2\bigr). \tag{8.4.70}
\end{equation}
对正数 $\lambda_+$，由
\begin{equation}
  |z_1|^2+|z_2|^2=2\lambda_++|z_3|^2+|z_4|^2>0 \tag{8.4.71}
\end{equation}
可知 $\mu_K^{-1}(\lambda_+)$ 的任意点 $(z_1,z_2,z_3,z_4)$ 都满足 $(z_1,z_2)\ne(0,0)$，因此可定义映射
\begin{equation}
  X_+:=\C^4/\!/_{\lambda_+}K\to\PP^1,\qquad [(z_1,z_2,z_3,z_4)]\mapsto[z_1:z_2]. \tag{8.4.72}
\end{equation}
把射影直线上的\emph{重言线丛}（tautological bundle）记作 $\OO_{\PP^1}(-1)$，则通过该映射，$X_+$ 作为复流形与 $\PP^1$ 上两个重言线丛的向量丛直和的全空间 $\OO_{\PP^1}(-1)\oplus\OO_{\PP^1}(-1)$ 同型。同样地，对负数 $\lambda_-$ 可定义映射
\begin{equation}
  X_-:=\C^4/\!/_{\lambda_-}K\to\PP^1,\qquad [(z_1,z_2,z_3,z_4)]\mapsto[z_3:z_4], \tag{8.4.73}
\end{equation}
通过该映射，$X_-$ 作为复流形与 $\PP^1$ 上两个重言线丛的向量丛直和的全空间 $\OO_{\PP^1}(-1)\oplus\OO_{\PP^1}(-1)$ 同型。

% [Tutor Note: 「同義直線束」是原书用词，日文通常作「普遍直線束」，指 tautological bundle；此处按中文惯例译作重言线丛。]

$X_+$ 与 $X_-$ 都是环面簇，其中都含有作为稠密开轨道的代数环面，但无法把环面的恒等映射延拓为 $X_+$ 与 $X_-$ 的同型。然而，$X_+$ 与 $X_-$ 作为复流形、也作为辛环面流形都同型。

% [原文 p.117]
因子列为
\begin{equation}
  0\to\R^3\xrightarrow{\ \left(\begin{smallmatrix}0&0&1\\1&1&1\\1&0&1\\0&1&1\end{smallmatrix}\right)\ }\R^4
  \xrightarrow{\ \left(\begin{smallmatrix}1\\)\ }\R\to 0 \tag{8.4.74}
\end{equation}
由该式给出，动量多面体由
\begin{align}
  m_3&\ge\lambda_1, \tag{8.4.75}\\
  m_1+m_2+m_3&\ge\lambda_2, \tag{8.4.76}\\
  m_1+m_3&\ge\lambda_3, \tag{8.4.77}\\
  m_2+m_3&\ge\lambda_4 \tag{8.4.78}
\end{align}
定义。其中
\begin{equation}
  \lambda=\lambda_1+\lambda_2-\lambda_3-\lambda_4 \tag{8.4.79}
\end{equation}
取
\begin{equation}
  (\lambda_1,\lambda_2,\lambda_3,\lambda_4)=(0,\lambda,0,0), \tag{8.4.80}
\end{equation}
则由
\begin{equation}
  m_3\ge 0,\qquad m_1+m_3\ge 0,\qquad m_2+m_3\ge 0 \tag{8.4.81}
\end{equation}
定义的部分集合是由 $(1,0,0)$、$(0,1,0)$ 及 $(-1,-1,1)$ 张成的三维多面锥，用
\begin{equation}
  m_1+m_2+m_3\ge\lambda \tag{8.4.82}
\end{equation}
切割它所得的多面体就是 $X_\lambda$ 的动量多面体。当 $\lambda>0$ 时它以 $(\lambda,0,0)$ 及 $(0,\lambda,0)$ 为顶点，当 $\lambda<0$ 时它以 $(0,0,0)$ 及 $(\lambda,\lambda,-\lambda)$ 为顶点。

当 $\lambda=0$ 时，$\mu_K^{-1}(0)$ 含有原点 $0\in\C^4$。该点是 $\mu_K$ 的临界值，在此点处 $K$ 的作用不自由。由此可知，$X_0=\mu_K^{-1}(0)/K$ 在 $[0]$ 处具有奇点。与该环面簇对应的扇形 $\Sigma$ 只有唯一一个三维锥
\begin{equation}
  \sigma=\sigma_1+\sigma_2+\sigma_3+\sigma_4, \tag{8.4.83}
\end{equation}
$\sigma^\vee\cap M$ 由
\[
  p_1=(1,0,0),\quad p_2=(0,1,0),\quad p_3=(-1,0,1),\quad p_4=(0,-1,1)
\]
生成，并具有关系式

% [原文 p.118]
\begin{equation}
  p_1+p_3=p_2+p_4. \tag{8.4.84}
\end{equation}
因此半群环为
\begin{equation}
  \C[\sigma^\vee\cap M]\cong\C[z_1,z_2,z_3,z_4]/(z_1z_3-z_2z_4). \tag{8.4.85}
\end{equation}
由于扇形 $\Sigma$ 只有唯一一个三维锥 $\sigma$，对应的环面簇为
\begin{equation}
  X_\Sigma=U_\sigma \tag{8.4.86}
\end{equation}
由 (8.4.85)，
\begin{align}
  U_\sigma&\cong\operatorname{Spec}\C[z_1,z_2,z_3,z_4]/(z_1z_3-z_2z_4) \tag{8.4.87}\\
  &\cong\{(z_1,z_2,z_3,z_4)\in\C^4\mid z_1z_3-z_2z_4=0\}. \tag{8.4.88}
\end{align}
它在 $(z_1,z_2,z_3,z_4)=(0,0,0,0)$ 处具有称为\emph{普通二重点}（ordinary double point）的奇点。经变量变换
\begin{align}
  z_1&=w_1+\sqrt{-1}w_2, & z_2&=-(w_3-\sqrt{-1}w_4), \tag{8.4.89}\\
  z_3&=w_1-\sqrt{-1}w_2, & z_4&=w_3-\sqrt{-1}w_4, \tag{8.4.90}
\end{align}
可知该奇点是 Morse 型（或称 $A_1$ 型）奇点；
\begin{equation}
  \{(w_1,w_2,w_3,w_4)\in\C^4\mid w_1^2+w_2^2+w_3^2+w_4^2=0\}. \tag{8.4.91}
\end{equation}
$X_+$ 与 $X_-$ 分别是 $X_0$ 沿理想 $(z_1,z_2)$ 与 $(z_3,z_4)$ 的爆破。反过来说，把 $X_+$ 或 $X_-$ 视为 $\PP^1$ 上两个重言线丛的直和，并把零截面的像收缩，便得到 $X_0$。

% [Tutor Note: Atiyah flop 是三维极小模型理论中最早的分裂收缩例子：$X_+\to X_0$ 与 $X_-\to X_0$ 沿同一条 $\PP^1$ 收缩，把两个爆破互相替换，得到复流形同型但环面结构不同的两族辛环面流形。]

\chapterreview{本章从辛几何与代数几何两个方向研究环面流形。辛几何方向以 Delzant 对应为核心：紧 $2n$ 维辛流形上 $n$ 维环面的有效 Hamilton 作用由其动量多面体（Delzant 多面体）完全决定，局部标准形与 Moser 定理给出唯一性。代数几何方向以扇形为组合数据，把仿射概形 $\operatorname{Spec}\C[\sigma^\vee\cap M]$ 粘贴成环面簇 $X_\Sigma$，而两者的桥梁是法扇形：动量多面体的法扇形给出与原辛环面流形微分同相的环面簇。§8.3 的反射多面体与 Fano 多面体把 Fano 环面簇的分类化为格多面体的分类。§8.4 用射影直线、射影平面与 Atiyah flop 三个例子，演示了从多面体数据复原复结构、以及辛商在临界值处产生奇点并可被两种方式爆破的全过程。}
